\documentclass[10pt,a4paper]{article}

\usepackage[T1]{fontenc}
\usepackage[utf8]{inputenc}
\usepackage[ngerman]{babel}
\usepackage[a4paper,margin=14mm]{geometry}
\usepackage{lmodern}
\usepackage{xcolor}
\usepackage{tabularx}
\usepackage{array}
\usepackage{listings}
\usepackage{enumitem}
\usepackage{microtype}

\pagestyle{empty}
\setlength{\parindent}{0pt}
\setlength{\parskip}{0pt}
\renewcommand{\familydefault}{\sfdefault}

\definecolor{Navy}{HTML}{17324D}
\definecolor{Blue}{HTML}{3B6EA5}
\definecolor{Sky}{HTML}{EAF3FA}
\definecolor{Mint}{HTML}{E8F5EF}
\definecolor{Gold}{HTML}{F3B33D}
\definecolor{Ink}{HTML}{1E2935}
\definecolor{Muted}{HTML}{617080}
\definecolor{Line}{HTML}{D8E1E8}
\definecolor{CodeBg}{HTML}{F5F7F9}

\color{Ink}

\newcommand{\sectiontitle}[2]{%
  \vspace{3mm}
  {\color{#1}\bfseries\large #2}\par
  \vspace{1mm}
  {\color{#1}\rule{\linewidth}{0.7pt}}\par
  \vspace{1.5mm}
}

\newcommand{\solutionrow}[3]{%
  \textcolor{Blue}{\bfseries #1} &
  \ttfamily\small #2 &
  \small #3 \\[2.3mm]
}

\lstset{
  basicstyle=\ttfamily\small,
  backgroundcolor=\color{CodeBg},
  frame=single,
  rulecolor=\color{Line},
  framesep=3mm,
  xleftmargin=1mm,
  xrightmargin=1mm,
  aboveskip=1.5mm,
  belowskip=1mm,
  showstringspaces=false,
  columns=fullflexible
}

\begin{document}

\colorbox{Navy}{%
  \parbox{\dimexpr\linewidth-2\fboxsep\relax}{%
    \vspace{3mm}
    \color{white}
    \begin{tabularx}{\linewidth}{@{}Xr@{}}
      {\small\bfseries INFORMATIK · KLASSE 9} &
      \colorbox{Gold}{\color{Navy}\small\bfseries\strut LÖSUNGEN}
    \end{tabularx}
    \vspace{2.5mm}

    {\fontsize{22}{25}\selectfont\bfseries Sicherungsblatt: Aufgabe 1}\par
    \vspace{1mm}
    {\large Programmtext analysieren}
    \vspace{3mm}
  }%
}

\vspace{3mm}
\colorbox{Sky}{%
  \parbox{\dimexpr\linewidth-2\fboxsep\relax}{%
    \vspace{1.8mm}
    \textcolor{Blue}{\bfseries Ziel:}
    Du kannst Objekte, Methoden und Parameter in einem kurzen Java-Programm
    erkennen und ihre Wirkung erklären.
    \vspace{1.8mm}
  }%
}

\sectiontitle{Blue}{1 · Das Ausgangsprogramm}

\begin{lstlisting}
Circle ball = new Circle(200, 50, 50);
ball.move(10, 10);
ball.setFillColor(Color.red);
ball.destroy();
\end{lstlisting}

\sectiontitle{Blue}{2 · Musterlösung zu b): Zeile für Zeile}

\renewcommand{\arraystretch}{1.05}
\begin{tabularx}{\linewidth}{@{}>{\centering\arraybackslash}p{7mm} >{\raggedright\arraybackslash}p{67mm} X@{}}
  \solutionrow{1}{Circle ball = new Circle(200, 50, 50);}{Ein neues Kreis-Objekt wird erzeugt und in der Objektvariablen \texttt{ball} gespeichert. Die Zahlen legen Position und Größe fest.}
  \solutionrow{2}{ball.move(10, 10);}{Der Kreis wird um die angegebenen Werte in x- und y-Richtung verschoben.}
  \solutionrow{3}{ball.setFillColor(Color.red);}{Die Füllfarbe des Kreises wird auf Rot gesetzt.}
  \solutionrow{4}{ball.destroy();}{Der Kreis wird wieder von der Zeichenfläche entfernt.}
\end{tabularx}

\sectiontitle{Blue}{3 · Musterlösung zu c): Was bedeuten die Zahlen?}

\begin{tabularx}{\linewidth}{@{}p{50mm} X@{}}
  \colorbox{Sky}{\parbox{46mm}{%
    \texttt{new Circle(200, 50, 50)}

    \vspace{1.5mm}
    \textcolor{Blue}{\bfseries 200} = x-Position

    \textcolor{Blue}{\bfseries 50} = y-Position

    \textcolor{Blue}{\bfseries 50} = Radius / Größe
  }} &
  \colorbox{Mint}{\parbox{\dimexpr\linewidth-54mm\relax}{%
    \texttt{move(10, 10)}

    \vspace{1.5mm}
    \textcolor{Blue}{\bfseries erste 10} = Verschiebung in x-Richtung

    \textcolor{Blue}{\bfseries zweite 10} = Verschiebung in y-Richtung

    \vspace{1mm}
    Der Kreis verändert dabei seine Größe und Farbe nicht.
  }}
\end{tabularx}

\vspace{2.5mm}
\colorbox{CodeBg}{%
  \parbox{\dimexpr\linewidth-2\fboxsep\relax}{%
    \vspace{1.5mm}
    \textcolor{Blue}{\bfseries Gute Versuchsstrategie:}
    Verändere immer nur eine Zahl. Starte das Programm erneut und vergleiche
    die Beobachtung mit dem vorherigen Durchlauf.
    \vspace{1.5mm}
  }%
}

\sectiontitle{Blue}{4 · Ablauf und Merksätze}

\begin{center}
  \colorbox{Sky}{\strut\bfseries Kreis erzeugen}
  \textcolor{Gold}{\bfseries\quad→\quad}
  \colorbox{Sky}{\strut\bfseries bewegen}
  \textcolor{Gold}{\bfseries\quad→\quad}
  \colorbox{Sky}{\strut\bfseries rot färben}
  \textcolor{Gold}{\bfseries\quad→\quad}
  \colorbox{Sky}{\strut\bfseries entfernen}
\end{center}

\vspace{1mm}
\begin{tabularx}{\linewidth}{@{}p{7mm}X@{}}
  \textcolor{Gold}{\Large\bfseries !} &
  \textbf{Objekt und Variable:} \texttt{new Circle(...)} erzeugt das Objekt;
  \texttt{ball} ist die Objektvariable, über die es angesprochen wird. \\[1.5mm]
  \textcolor{Gold}{\Large\bfseries !} &
  \textbf{Methodenaufruf:} Mit der Punktnotation, zum Beispiel
  \texttt{ball.move(...)}, wird eine Methode für genau dieses Objekt aufgerufen. \\[1.5mm]
  \textcolor{Gold}{\Large\bfseries !} &
  \textbf{Parameter:} Werte in den Klammern steuern, wie ein Objekt erzeugt
  oder wie eine Methode ausgeführt wird.
\end{tabularx}

\vfill
{\color{Line}\rule{\linewidth}{0.5pt}}\par
\vspace{1.5mm}
\begin{tabularx}{\linewidth}{@{}Xr@{}}
  {\small\color{Muted}Klasse 9 · Modellierung und Programmierung} &
  {\small\color{Muted}Sicherungsblatt · Aufgabe 1}
\end{tabularx}

\end{document}
